% -*- coding: utf-8 -*-
% Adapted from the user-supplied LaTeX模板.zip, main.tex (designed by zcf).
% Preserved: NKU cover, logo, margins, Chinese numbering, contents, headers, watermark.
\documentclass[UTF8,a4paper,10pt,fontset=fandol]{ctexart}
\usepackage[left=3.17cm,right=3.17cm,top=2.74cm,bottom=2.74cm]{geometry}
\IfFileExists{NotoSerifCJK-Regular.ttc}{
  \setCJKmainfont[FontIndex=2,BoldFont=NotoSerifCJK-Bold.ttc,
    ItalicFont=NotoSansCJK-Regular.ttc]{NotoSerifCJK-Regular.ttc}
  \setCJKsansfont[FontIndex=2,BoldFont=NotoSansCJK-Bold.ttc]{NotoSansCJK-Regular.ttc}
}{
  \IfFontExistsTF{Noto Serif CJK SC}{\setCJKmainfont{Noto Serif CJK SC}}{
    \IfFileExists{C:/Windows/Fonts/simsun.ttc}{
      \setCJKmainfont[FontIndex=0,AutoFakeBold=2]{C:/Windows/Fonts/simsun.ttc}
    }{}
  }
}
\usepackage{amsmath,amssymb,graphicx,booktabs,array,tabularx,enumitem}
\usepackage{caption,float,xcolor,listings,xurl,fancyhdr,tikz,eso-pic,needspace}
\usepackage[colorlinks,linkcolor=black,citecolor=black,urlcolor=blue!55!black]{hyperref}
\hypersetup{pdftitle={编译流程与多层中间表示实验报告},pdfauthor={柯云超；郑奕杰}}
\graphicspath{{lab0/report/assets/}{assets/}}
\ctexset{
  section={name={,、},number=\chinese{section},format=\Large\bfseries\centering},
  subsection={name={（,）},number=\chinese{subsection},format=\large\bfseries},
  subsubsection={name={,.},number=\arabic{subsubsection},format=\normalsize\bfseries}
}
\pagestyle{fancy}
\fancyhf{}
\fancyhead[L]{\small\kaishu\leftmark}
\fancyhead[R]{\small\kaishu 编译器课程预备实验报告}
\fancyfoot[C]{\thepage}
\setlength{\headheight}{15pt}
\renewcommand{\headrulewidth}{0.1pt}
\renewcommand{\footrulewidth}{0pt}
\renewcommand{\sectionmark}[1]{\markboth{\chinese{section}、#1}{}}
\newcommand{\HRule}{\rule{\linewidth}{0.5mm}}
\newcommand{\sihao}{\fontsize{14pt}{21pt}\selectfont}
\newcommand{\watermark}{\AddToShipoutPictureBG{
  \begin{tikzpicture}[remember picture,overlay]
  \node[rotate=60,scale=10,text=gray!9] at (current page.center) {NKU};
  \end{tikzpicture}}}
\lstset{basicstyle=\ttfamily\small,breaklines=true,columns=fullflexible,
  keepspaces=true,showstringspaces=false,numbers=left,numberstyle=\tiny\color{blue!50},
  numbersep=6pt,xleftmargin=1.5em,frame=trbl,rulecolor=\color{black!25},
  keywordstyle=\color{blue!70!black},commentstyle=\color{black!55},
  aboveskip=0.7em,belowskip=0.7em}
\captionsetup{font=small,labelfont=bf}
\setlength{\emergencystretch}{3em}
\setlength{\parskip}{0.25em}
\setlist{nosep,leftmargin=2em}
\renewcommand{\arraystretch}{1.18}

\begin{document}
\hypersetup{pageanchor=false}
\begin{titlepage}
\begin{center}
  \includegraphics[width=0.8\textwidth]{NKU.png}\\[0.8cm]
  {\Huge\kaishu\bfseries 南\quad 开\quad 大\quad 学}\\[0.65cm]
  {\huge\kaishu\bfseries 计\ 算\ 机\ 学\ 院}\\[0.45cm]
  {\Large\bfseries 编译器课程预备工作实验报告}\\[0.6cm]
  \HRule\\[0.65cm]
  {\LARGE\bfseries 编译流程与多层中间表示}\\[0.3cm]
  {\Large SysY、LLVM IR、RISC-V 与 MLIR}\\[0.55cm]
  \HRule\\[0.85cm]
  {\Large\kaishu 柯云超\quad 2413575}\\[0.35cm]
  {\Large\kaishu 郑奕杰\quad 2412677}\\[0.25cm]
  {\small 代码仓库：\url{https://github.com/kyc001/compiler}}\\[0.4cm]
\end{center}
\noindent\textbf{小组分工安排：}
\begin{itemize}
\item \textbf{柯云超：}负责实验总体方案与技术路线设计，主导 SysY、LLVM IR 和 RISC-V 汇编的核心实现与语义等价验证；完成 SSA 与优化效果分析、MLIR 逐层降低实验及 AscendNPU IR 调研，构建自动化验证流程，统筹报告的技术分析与结论论证。
\item \textbf{郑奕杰：}负责阶乘样例的预处理、编译、汇编与链接基础实验，整理各阶段命令及输出；执行基础与边界用例，核对预期结果，协助整理实验环境说明、结果图表和参考资料。
\item 共同工作：交叉检查实验代码、核对结果与引用，复核并提交同一份报告。
\end{itemize}
\vfill
\begin{center}{\Large 2026年10月10日}\end{center}
\end{titlepage}

\hypersetup{pageanchor=true}
\pagenumbering{Roman}
\thispagestyle{plain}
\renewcommand{\abstractname}{\kaishu\sihao\textbf{摘要}}
\begin{abstract}
本实验以阶乘程序研究预处理、编译、汇编和链接的完整过程，观察宏替换、前端诊断、SSA 提升、机器指令、符号及重定位，并比较优化等级、调试选项与静态链接的影响。独立编写三组 SysY、LLVM IR 与 RISC-V64 等价程序，覆盖短路求值、作用域、数组、递归、整数除余及浮点运算与转换；34 个输入在五条执行路径上完成 170 次比对，全部通过，阶乘的五种构建结果另通过 85 次检查。进阶实验实测通用 MLIR 向量加的逐层降低，以三组输入／输出布局验证 12 个元素，并结合官方 VecAdd 分析 AscendNPU IR 的方言职责与数据依赖。结果表明，低层表示必须显式维护控制流、数据布局与调用约定，文件体积不能直接代表性能。Ascend 专用部分采用文档调研方式，未执行 BishengIR/CANN 或 NPU 程序。

\medskip
\noindent\textbf{关键词：}编译流程；SysY；LLVM IR；RISC-V；MLIR
\end{abstract}
\begingroup
\small
\setlength{\parskip}{0pt}
\tableofcontents
\endgroup
\clearpage
\pagenumbering{arabic}
\watermark

\section{引言}
编译器将具有类型、作用域和结构化控制流的源程序转换为机器可执行代码。仅观察最终输出，难以区分前端检查、中端优化、后端生成以及汇编链接各自承担的工作。本实验围绕课程预备要求\cite{assignment}，以一个阶乘程序研究完整语言处理过程，再以三组 SysY 程序验证手写 LLVM IR 和 RISC-V 汇编的行为，最后探索 MLIR 多级表示。

实验按“源程序或编译选项发生变化，阶段输出如何变化”的方法组织。基础概念结合实际输出说明，重点分析短路副作用、SSA、数组地址计算、浮点转换及多级降低。报告中的数据来自同一批自动化实验；保存的结果、命令和关键输出位于 \path{lab0/report/evidence/}，可通过 \texttt{make prep} 复现。测试验证给定输入上的一致性，不构成任意程序上的等价性证明。

\section{实验环境与设计}
\subsection{工具链和运行环境}
实验在 Ubuntu 22.04 的 WSL 环境完成，宿主架构为 x86\_64。表\ref{tab:tools}列出实际使用的工具。RISC-V 程序以 QEMU 用户模式执行，使用 RV64GC 指令集和 LP64D ABI。
\begin{table}[H]
\centering\small
\caption{实际使用的主要工具}\label{tab:tools}
\begin{tabularx}{\textwidth}{lX}
\toprule
工具或配置 & 版本或参数\\
\midrule
GCC & Ubuntu GCC 12.3.0\\
Clang / LLVM / MLIR & 15.0.7\\
RISC-V GCC & SiFive GCC-Metal 10.2.0\\
QEMU 用户模式 & 6.2.0\\
RISC-V 链接 & \texttt{-static -mcmodel=medany}\\
地址与重定位选项 & \texttt{-Wl,--no-relax,-Ttext=0x90000000}\\
\bottomrule
\end{tabularx}
\end{table}

运行库来自课程提供的 \path{lab0/runtime/lib.tar.gz}\cite{runtime,guide}。宿主程序链接其中的 \texttt{libsysy\_x86.a}。预编译 RISC-V 库的 ISA 属性与本机 GCC 10 链接器不兼容，因此按讲义允许的方法，从未经修改的 \texttt{sylib.c} 重建兼容库。课程头文件定义了计时全局变量，C 参考程序采用 \texttt{-fcommon} 编译；标准错误上的运行库计时信息单独记录，功能输出按标准输出比较。

\subsection{实验对象与验证口径}
阶乘程序用于编译流程实验，合法输入为 $0\le n\le12$；$12!=479001600$ 仍在 32 位有符号整数范围内，$13!$ 越界，因此先检查范围。等价编程使用 \texttt{features}、\texttt{arrays\_float} 和 \texttt{numeric\_edges}，分别测试控制流、二维浮点数组及数值边界。通过强制包含课程头文件补充运行库声明，建立 C 编译器可处理的参考路径；它不依赖尚未实现完整功能的课程编译器。SysY 浮点字面量为 float，C 中无后缀字面量为 double；本实验选取二进制可精确表示的常量及有界输入，各浮点中间结果也可精确表示，使这一差异不影响所测结果，不能由此断言二者对任意源码兼容。

预期值保存在独立用例文件中。每次执行同时检查退出码和标准输出，保留错误诊断；编译及运行失败会中止脚本。IR 先经过 LLVM verifier，随后才进入目标生成和执行。文件大小、函数符号大小与测试次数采用不同口径，避免将调试信息体积误当作指令数量。

\section{完整编译流程实验}
\subsection{程序与阶段输出}
主程序用 \texttt{scanf} 读入并检查整数，调用阶乘函数后用 \texttt{printf} 输出。其核心代码如下：
\par\noindent\begin{minipage}{\linewidth}
\begin{lstlisting}[language=C]
#ifndef FIRST_FACTOR
#define FIRST_FACTOR 2
#endif
int factorial(int n) {
    int result = 1;
    int i = FIRST_FACTOR;
    while (i <= n) {
        result = result * i;
        i = i + 1;
    }
    return result;
}
\end{lstlisting}
\end{minipage}\par

图\ref{fig:pipeline}说明处理过程，表\ref{tab:stages}对应可检查的产物。图中编译器内部各步骤是功能划分，具体工具可能将多个步骤融合实现。
\begin{figure}[H]
\centering
\fbox{\begin{minipage}{0.95\textwidth}
\centering\small
源程序 \texttt{.c}
$\xrightarrow{\text{预处理}}$ \texttt{.i}
$\xrightarrow{\text{编译}}$ \texttt{.s}
$\xrightarrow{\text{汇编}}$ \texttt{.o}
$\xrightarrow{\text{链接}}$ 可执行文件\\[0.65em]
编译器内部：Token $\rightarrow$ AST／语义检查 $\rightarrow$ IR／优化
$\rightarrow$ 指令选择、寄存器分配及汇编输出
\end{minipage}}
\caption{语言处理系统及编译器内部的功能分工}\label{fig:pipeline}
\end{figure}

\begin{table}[H]
\centering\small
\caption{阶段、观察命令及对应信息}\label{tab:stages}
\begin{tabularx}{\textwidth}{>{\raggedright\arraybackslash}p{1.2cm}>{\raggedright\arraybackslash}p{5.0cm}X}
\toprule
阶段 & 主要选项 & 观察内容\\
\midrule
预处理 & \texttt{gcc -E} & 宏与头文件展开、来源行标记\\
词法 & \texttt{clang-15 -Xclang -dump-tokens} & Token 类型、文本及位置\\
语法／语义 & \texttt{-Xclang -ast-dump}、\texttt{-fsyntax-only} & AST、类型与未声明标识符诊断\\
中间代码 & \texttt{-fdump-tree-gimple}、\texttt{-emit-llvm} & 显式分支、基本块与数据操作\\
优化／后端 & \texttt{opt-15 -passes=mem2reg}、\texttt{gcc -S} & SSA、指令和寄存器使用\\
汇编 & \texttt{gcc -c}、\texttt{objdump -dr} & 机器码与重定位位置\\
链接 & \texttt{gcc factorial.o -o factorial} & 符号解析、节布局与库依赖\\
\bottomrule
\end{tabularx}
\end{table}

\Needspace{8\baselineskip}
以下命令概括基本处理链；自动化脚本使用完整路径，并保留其参数\cite{gcc}。
\par\noindent\begin{minipage}{\linewidth}
\begin{lstlisting}[numbers=none]
gcc -E factorial.c -o factorial.i
gcc -S -O0 -g factorial.i -o factorial_O0.s
gcc -c factorial_O0.s -o factorial.o
gcc factorial.o -o factorial
\end{lstlisting}
\end{minipage}\par

\subsection{预处理、前端分析与语义错误}
预处理器展开 \texttt{\#include} 和宏、处理条件编译，消除注释并输出来源行标记。本实验中 \texttt{FIRST\_FACTOR} 在 \texttt{.i} 中被替换为 2；添加 \texttt{-DFIRST\_FACTOR=3} 后变为 3。输入 5 时，默认程序计算 $2\times3\times4\times5=120$，修改后的程序输出 $3\times4\times5=60$。宏改变了后续编译器看到的程序，预处理器本身不负责类型检查或生成机器码。

词法分析把 \texttt{while (i <= n)} 分解为关键字、括号、标识符和关系运算符等 Token，语法分析再建立 \texttt{WhileStmt}、条件表达式和循环体等 AST 节点。语义分析进一步完成名称绑定、类型检查和隐式转换。单独编译“返回未声明的 \texttt{missing\_name}”的函数时，Clang 报告 \texttt{use of undeclared identifier 'missing\_name'}，退出码为 1。程序的括号和语句结构合法仍不代表可以进入代码生成。

GCC 的 GIMPLE 输出把 \texttt{while} 改写为标签和跳转。实际转储中的核心语句为：
\par\noindent\begin{minipage}{\linewidth}
\begin{lstlisting}[numbers=none]
result = 1; i = 2;
goto <D.2386>;
<D.2387>:
result = result * i;
i = i + 1;
<D.2386>:
if (i <= n) goto <D.2387>; else goto <D.2385>;
\end{lstlisting}
\end{minipage}\par
源语言的循环嵌套在此变为控制流关系，后续分析便可围绕基本块和数据依赖进行。

\subsection{LLVM IR 与 SSA 提升}
使用 \texttt{-O0 -Xclang -disable-O0-optnone} 生成 LLVM IR，再单独执行 \texttt{mem2reg}，以隔离这一变换。禁用 \texttt{optnone} 是为了允许优化 pass 处理 O0 函数\cite{llvm}。
\par\noindent\begin{minipage}{\linewidth}
\begin{lstlisting}[numbers=none]
clang-15 -S -emit-llvm -O0 -Xclang -disable-O0-optnone \
  factorial.c -o factorial.raw.ll
opt-15 -S -passes=mem2reg factorial.raw.ll \
  -o factorial.ssa.ll
\end{lstlisting}
\end{minipage}\par
整个模块的 \texttt{alloca} 从 5 条降为 1 条，出现 3 条 \texttt{phi}。其中两条处理阶乘循环的累乘值和循环变量，另一条合并 main 的返回值。实际优化结果中的循环头如下，数值标签和变量名与转储一致：
\par\noindent\begin{minipage}{\linewidth}
\begin{lstlisting}[numbers=none]
2:  ; predecessors: %4, %1
  %.01 = phi i32 [ 1, %1 ], [ %5, %4 ]
  %.0 = phi i32 [ 2, %1 ], [ %6, %4 ]
  %3 = icmp sle i32 %.0, %0
  br i1 %3, label %4, label %7
\end{lstlisting}
\end{minipage}\par
\texttt{phi} 根据控制流前驱选择初始值或上轮结果，使每个 SSA 名只定义一次。main 的输入变量要将地址传给 \texttt{scanf}，无法由本次 pass 提升，故保留一条 \texttt{alloca}。这里的 IR 变量并非已经分配好的物理寄存器，后端仍须选择指令并决定哪些值驻留寄存器、哪些值溢出到栈。

\subsection{优化等级与调试选项的对照}
在 O0 的阶乘函数中，\texttt{result}、\texttt{i} 和 \texttt{n} 分别经 \texttt{-8(\%rbp)}、\texttt{-4(\%rbp)} 和 \texttt{-20(\%rbp)} 访问。O2 的核心循环如下：
\par\noindent\begin{minipage}{\linewidth}
\begin{lstlisting}[numbers=none]
.L3:
  imull %eax, %edx
  addl  $1, %eax
  cmpl  %edi, %eax
  jne   .L3
\end{lstlisting}
\end{minipage}\par
O2 将循环值保存在寄存器中，并将边界整理为与 $n+1$ 比较；调试指令 \texttt{.loc} 仍用于关联源代码。用 \texttt{nm -S} 读取函数符号大小，用文件元数据读取整个可执行文件大小，结果见表\ref{tab:options}。

\begin{table}[H]
\centering\small
\caption{同一源码的选项对照；大小单位为字节}\label{tab:options}
\begin{tabular}{lrrc}
\toprule
构建方式 & 可执行文件大小 & factorial 符号大小 & 含 \texttt{.debug\_info}\\
\midrule
O0、带 \texttt{-g} & 17\,704 & 54 & 是\\
O0、不带 \texttt{-g} & 16\,104 & --- & 否\\
O2、带 \texttt{-g} & 18\,768 & 56 & 是\\
O0、带 \texttt{-g}、静态链接 & 1\,035\,896 & --- & 是\\
\bottomrule
\end{tabular}
\end{table}
O2 改善了循环的寄存器使用，但其函数大小略增，程序文件也更大；指令布局、对齐与调试信息都会影响体积。因此，不能从文件更小推断程序更快，也不能从 O2 名称推断本实验测得了加速。本实验未进行可信的性能计时，结论限于实际观察到的结构变化。

O0、O2、静态链接、原始 LLVM IR 和 mem2reg 后 IR 五种程序，分别验证 $0,\ldots,12$、$-1$、13、非数字输入与空输入，共 $17\times5=85$ 次检查，全部通过。非法输入无标准输出且返回 1，合法输入输出对应阶乘。另行验证宏修改后的结果为 60，避免把有意改变语义的变体纳入等价结论。

\subsection{汇编、符号与链接}
汇编器把目标相关汇编转换为机器码，建立节、符号表和重定位记录。\texttt{readelf -h} 显示 \texttt{factorial.o} 为 REL 类型；\texttt{nm -u} 显示尚未解析的 \texttt{printf}、\texttt{\_\_isoc99\_scanf} 和 \texttt{\_\_stack\_chk\_fail}。最后一个符号源于本机 GCC 默认栈保护，与源码显式调用无关。

\texttt{readelf -r} 中，\texttt{.rela.text} 内偏移 \texttt{0x68} 和 \texttt{0xa3} 分别存在指向输入和输出函数的 \texttt{R\_X86\_64\_PLT32} 重定位。汇编器已知道指令形式，但外部符号最终地址仍需链接器和动态装载机制处理。链接器组合目标文件、启动代码及库，解析符号、安排节并修正引用。默认链接程序带有动态解释器段；静态程序将所需库代码纳入文件，体积约为默认程序的 58.5 倍。两者在本实验输入上的输出一致，但依赖方式和产物布局不同。

\section{SysY、LLVM IR 与 RISC-V 等价编程}
\subsection{样例设计与语言特性}
按 SysY 定义及文法补充\cite{sysy}，独立编写三组源程序、LLVM IR 和 RISC-V 汇编。手写 IR 与手写汇编保存在 \path{lab0/prep/02_equivalence/}，与工具自动生成的文件分开，便于检查表示之间的对应关系。表\ref{tab:coverage}列出覆盖情况。
\begin{table}[H]
\centering\small
\caption{样例功能与语义覆盖}\label{tab:coverage}
\begin{tabularx}{\textwidth}{p{2.3cm}X p{3.1cm}}
\toprule
函数／样例 & 覆盖特性 & 可观察结果\\
\midrule
factorial & while、赋值、乘法、比较、返回 & $n!$\\
arithmetic & 单目正负、加乘除模 & $n+(n\bmod3)$\\
control & if、break、continue & 跳过 3，遇 9 停止累加\\
scope & 嵌套块、常量、同名变量隐藏 & $2n+7$\\
logic & 逻辑与或非、短路、全局副作用 & 合法输入时为 111\\
arrays\_float & 二维数组、数组形参、浮点加法、隐式转换 & 行求和、截断与混合运算\\
numeric\_edges & 负数除余、浮点四则与比较、递归、void 函数、局部数组补零与写入 & 商余、舍入、递归和、数组校验值\\
\bottomrule
\end{tabularx}
\end{table}

\subsection{控制流、短路与作用域}
\texttt{logic} 用 \texttt{bump()} 对全局计数器加一，使求值次数可以观察：
\par\noindent\begin{minipage}{\linewidth}
\begin{lstlisting}[language=C]
if ((n > 0 && bump() != 0) || bump() != 0)
    calls = calls + 10;
if (!(n < 0) && (n >= 0 || bump() != 0))
    calls = calls + 100;
\end{lstlisting}
\end{minipage}\par
对正数，第一次条件只执行左侧 \texttt{bump}；对 0，逻辑与跳过左侧调用，再执行逻辑或右侧。第二个条件在合法输入上不执行 \texttt{bump}，因此结果为 $1+10+100=111$。LLVM IR 通过条件 \texttt{br} 把函数调用放在对应路径中，汇编采用条件跳转实现同一控制流；若无条件计算逻辑运算的两个操作数，计数器就会暴露错误。

循环中的 \texttt{continue} 跳向下一次条件判断，\texttt{break} 跳到出口。LLVM 的 \texttt{phi} 在循环头合并不同前驱的循环变量和累计值。块作用域的同名变量则在名称解析后对应不同值：内层 \texttt{x=7} 不覆盖外层 \texttt{x=n} 的存储含义，故 \texttt{scope(n)} 返回 $2n+7$。词法作用域与基本块属于不同层次，不能用基本块边界替代作用域管理。

\Needspace{20\baselineskip}
\subsection{整数语义与 RISC-V 调用约定}
SysY 整数使用 32 位表示，手写 IR 采用 \texttt{i32}、\texttt{sdiv} 和 \texttt{srem}；RV64 汇编使用 \texttt{addw}、\texttt{mulw}、\texttt{divw} 和 \texttt{remw} 保持相应整数运算宽度。阶乘在手写汇编中的核心为：
\par\noindent\begin{minipage}{\linewidth}
\begin{lstlisting}[numbers=none]
    li t0, 1
    li t1, 1
.Lfactorial_loop:
    bgt t1, a0, .Lfactorial_exit
    mulw t0, t0, t1
    addiw t1, t1, 1
    j .Lfactorial_loop
.Lfactorial_exit:
    mv a0, t0
    ret
\end{lstlisting}
\end{minipage}\par
整数参数和返回值使用 \texttt{a0} 等参数寄存器；跨调用要保存的值放入 \texttt{s0} 等被调用者保存寄存器。非叶函数保存 \texttt{ra} 与实际使用的保存寄存器，栈保持 16 字节对齐。浮点参数和返回值遵循 LP64D，使用 \texttt{fa0} 等寄存器\cite{riscv}。正确的局部指令序列若破坏调用约定，仍可能在运行库调用后出现错误。

\subsection{数组地址与浮点转换}
数组样例定义 \texttt{weights[2][2]=\{\{1.5,2.0\},\{3.25,4.0\}\}}，通过形参 \texttt{float values[][2]} 传入行求和函数。元素 \texttt{values[row][col]} 的字节偏移为
\[
  \operatorname{offset}=4(2\,row+col),
\]
因为每个单精度元素占 4 字节、每行两个元素。LLVM IR 用 \texttt{getelementptr} 按类型计算地址；手写汇编将行号乘 8，再用 \texttt{flw} 读取两列并以 \texttt{fadd.s} 相加。

浮点和赋给整数时，LLVM 使用 \texttt{fptosi}，汇编使用带向零舍入模式的 \texttt{fcvt.w.s ..., rtz}。整数行号参与浮点加法时，分别使用 \texttt{sitofp} 和 \texttt{fcvt.s.w}。两行的和为 3.5 与 7.25，截断后为 3 与 7；加入行号后输出浮点值 3.5 与 8.25。课程运行库的 \texttt{putfloat} 以十六进制浮点格式打印，分别对应 \texttt{0x1.cp+1} 与 \texttt{0x1.08p+3}。这种输出格式是运行库约定，不是数值计算错误。

以下为手写 \texttt{row\_sum} 的 LLVM IR 主体；它对应源码的 \texttt{values[row][0]+values[row][1]}，没有省略实际的数组读取：
\par\noindent\begin{minipage}{\linewidth}
\begin{lstlisting}[numbers=none]
%index = sext i32 %row to i64
%left.ptr = getelementptr [2 x float], ptr %values, i64 %index, i64 0
%right.ptr = getelementptr [2 x float], ptr %values, i64 %index, i64 1
%left = load float, ptr %left.ptr, align 4
%right = load float, ptr %right.ptr, align 4
%sum = fadd float %left, %right
ret float %sum
\end{lstlisting}
\end{minipage}\par
\texttt{getelementptr} 只计算地址，\texttt{load} 才读取元素。形参首维退化为指针，故 GEP 的源元素类型是“一行”的 \texttt{[2 x float]}。汇编用 \texttt{slli t0,a1,3} 计算行偏移，两个 \texttt{flw} 访问偏移 0 和 4，\texttt{fadd.s} 将结果放入 \texttt{fa0} 返回；三种表示因此具有相同的行跨度、元素宽度和返回类型。

\subsection{负数运算、递归与数组写入的补充验证}
只使用非负数不能区分向零取整与向下取整。\texttt{numeric\_edges} 读取 $n,d$，限制 $-12\le n\le12$、$1\le|d|\le3$，在除法之前拒绝零除数和越界值。按语言规则，$q=\operatorname{trunc}(n/d)$、$r=n-qd$；因此 $-7/3=-2$、$-7\bmod3=-1$，而 $7/(-3)=-2$、$7\bmod(-3)=1$。若测试预期直接采用 Python 的 \texttt{//} 和 \texttt{\%}，会误用向下取整规则。本实验在独立 JSON 中显式保存这些预期值。

浮点函数计算 $y=(((-x/2)-0.25)+1)\times2/4$，其中 $x$ 由 $n$ 转换而来。手写 IR 依次使用 \texttt{fneg/fdiv/fsub/fadd/fmul/fdiv}，汇编使用对应单精度指令；$n=7$ 时 $y=-1.375$，转为整数应为 $-1$。另外以浮点小于比较输出符号标记。递归函数计算 $S(k)=k+S(k-1)$、$S(0)=0$，令 $k=|n|$，检查多层调用时返回地址和保存寄存器是否保持正确。

局部数组声明为 \texttt{int values[3]=\{n\}}，后两个元素必须补零；\texttt{void update(int values[],int delta)} 通过数组形参读写后，数组应为 $(n,d,n-d)$。输出加权校验值 $v_0+2v_1+3v_2=4n-d$，使缺失初始化或错误写入更易被观察。该样例共有 15 个输入，包括商余的四种符号组合、递归基例、正负范围端点和五个非法输入。

\subsection{链接路径及结果}
五条路径在每个用例上分别执行，结果见表\ref{tab:results}。其中“源程序”路径将本实验 SysY 子集按 C 语法编译；“独立 IR”路径使用手写 \texttt{.ll}；“手写汇编”路径直接汇编链接。
\begin{table}[H]
\centering\small
\caption{五条执行路径的等价比对结果}\label{tab:results}
\begin{tabularx}{\textwidth}{Xcccc}
\toprule
执行路径 & 控制流 & 数组 & 数值边界 & 合计\\
\midrule
源程序 $\rightarrow$ Clang $\rightarrow$ 宿主程序 & 15/15 & 4/4 & 15/15 & 34/34\\
独立 IR $\rightarrow$ Clang $\rightarrow$ 宿主程序 & 15/15 & 4/4 & 15/15 & 34/34\\
手写汇编 $\rightarrow$ 交叉链接 $\rightarrow$ QEMU & 15/15 & 4/4 & 15/15 & 34/34\\
源程序 $\rightarrow$ RISC-V GCC $\rightarrow$ QEMU & 15/15 & 4/4 & 15/15 & 34/34\\
独立 IR $\rightarrow$ llc $\rightarrow$ 链接 $\rightarrow$ QEMU & 15/15 & 4/4 & 15/15 & 34/34\\
\midrule
全部 & 75/75 & 20/20 & 75/75 & 170/170\\
\bottomrule
\end{tabularx}
\end{table}
\texttt{features} 输入 $-1,0,\ldots,12,13$；\texttt{arrays\_float} 输入 $-1,0,1,2$。越界输入输出 \texttt{-1}，且在乘法或数组访存之前退出。典型输出见表\ref{tab:examples}，每个单元格对应一行标准输出。
\begin{table}[H]
\centering\small
\caption{代表性输入与五条路径一致的输出}\label{tab:examples}
\begin{tabular}{lll}
\toprule
样例 & 输入 & 输出\\
\midrule
features & 0 & \texttt{1 0 0 7 111}\\
features & 5 & \texttt{120 7 12 17 111}\\
features & 12 & \texttt{479001600 12 33 31 111}\\
features & 13 & \texttt{-1}\\
arrays\_float & 0 & \texttt{0x1.cp+1 3}\\
arrays\_float & 1 & \texttt{0x1.08p+3 7}\\
arrays\_float & 2 & \texttt{-1}\\
\bottomrule
\end{tabular}
\end{table}
数值边界样例按“商、余数、浮点结果、截断值、负数标记、递归和、数组校验值”输出。例如输入 \texttt{7 -3} 得到 \texttt{-2 1 -0x1.6p+0 -1 1 28 31}。所有路径都链接课程 SysY 运行库，退出码均为 0；非法输入按样例约定输出 \texttt{-1}。测试覆盖三组程序的正常路径和边界，尚未穷举全部 SysY 语法、NaN／无穷值或任意浮点输入；本实验也不将整数溢出、除零或越界访存作为合法算术结果比较。

\section{MLIR 与 AscendNPU IR 探索}
\subsection{通用 VecAdd 的逐层降低}
MLIR 允许在同一框架中保留不同抽象层级的操作和类型，并通过转换 pass 逐步引入目标约束\cite{mlir}。本地设计四元素单精度向量加，起始表示含 \texttt{scf.for}、\texttt{arith.addf} 及 \texttt{memref.load/store}。输入
\[
a=(1,2,3,4),\qquad b=(0.5,2.5,4.5,6.5),
\]
期望输出为 $(1.5,4.5,7.5,10.5)$。

\texttt{convert-scf-to-cf} 将结构化循环转换为显式控制流，实际输出的循环头如下：
\par\noindent\begin{minipage}{\linewidth}
\begin{lstlisting}[numbers=none]
cf.br ^bb1(%c0 : index)
^bb1(%0: index):
  %1 = arith.cmpi slt, %0, %c4 : index
  cf.cond_br %1, ^bb2, ^bb3
\end{lstlisting}
\end{minipage}\par
循环变量成为块参数，回边通过 \texttt{cf.br} 传入更新值；其数据流作用对应 LLVM 中 \texttt{phi} 的合并机制。后续转换将 MemRef、Arith、Func 和 CF 依次降至 LLVM 方言，并消除转换过程中遗留的类型桥接操作。核心命令为：
\par\noindent\begin{minipage}{\linewidth}
\begin{lstlisting}[numbers=none]
mlir-opt-15 01_structured.mlir --convert-scf-to-cf \
  -o 02_cfg.mlir
mlir-opt-15 02_cfg.mlir --convert-memref-to-llvm \
  --convert-arith-to-llvm --convert-func-to-llvm \
  --convert-cf-to-llvm --reconcile-unrealized-casts \
  -o 03_llvm_dialect.mlir
mlir-translate-15 --mlir-to-llvmir \
  03_llvm_dialect.mlir -o 04_vecadd.ll
\end{lstlisting}
\end{minipage}\par
MemRef 降低后需要显式表达指针、偏移、维度和步长。实验通过 \texttt{llvm.emit\_c\_interface} 生成 C 接口，以包含 allocated/aligned 指针、offset、sizes 和 strides 的一维描述符传入数组。这里的静态类型 \texttt{memref<4xf32>} 要求 offset 为 0、size 为 4、stride 为 1；描述符含有这些字段并不表示该函数支持任意步长。三个 MemRef 在内部函数签名中展开为 15 个参数，外层 C 包装函数仍接收三个描述符指针。

表\ref{tab:lowering}以保存的转储对照同一循环。LLVM 方言的 \texttt{llvm.fadd} 仍是 MLIR 操作；经 \texttt{mlir-translate} 后才成为 LLVM IR 的 \texttt{fadd} 指令，两者不能只因名称相近便混为同一表示。
\begin{table}[H]
\centering\small
\caption{VecAdd 核心循环在各层的实测变化}\label{tab:lowering}
\begin{tabularx}{\textwidth}{lXX}
\toprule
阶段 & 控制流与数据表示 & 数值／内存操作\\
\midrule
结构化 MLIR & \texttt{scf.for}、\texttt{index} 循环变量 & 2 次 \texttt{memref.load}、1 次加法、1 次写回\\
CF 层 & \texttt{cf.cond\_br} 和块参数 & MemRef 访存及 \texttt{arith.addf} 保留\\
LLVM 方言 & \texttt{llvm.cond\_br}、\texttt{i64} 索引、描述符 & 3 个 GEP、2 个 load、1 个 \texttt{llvm.fadd}、1 个 store\\
LLVM IR & \texttt{br}、\texttt{phi i64} 和指针 & \texttt{getelementptr/load/fadd/store}\\
\bottomrule
\end{tabularx}
\end{table}
最终 IR 中 \texttt{\%32 = phi i64 [ \%44, \%34 ], [ 0, \%15 ]} 接收初值和回边值；循环体的 \texttt{\%41 = fadd float \%37, \%40} 对应起始的 \texttt{arith.addf}。加法本身保持一一对应，新增操作主要承担控制流、地址计算和 ABI 适配，故降低过程不能仅按 IR 行数评价优劣。

最终 LLVM IR 经 verifier 检查，再由 Clang 与 C 驱动链接。三组逐元素检查分别得到：基本输入 $(1.5,4.5,7.5,10.5)$；正负混合输入 $(-1,-2,0,-0.125)$；将输出与第一输入设为同一数组时仍得到后一结果，共 $3\times4=12$ 个元素通过。原地测试只覆盖两者完全重合，不扩展为任意部分重叠的保证。各数值均可由二进制精确表示，故此处使用等值比较。

\subsection{官方 Ascend VecAdd 与方言分工}
课程进阶要求允许在环境配置困难时通过阅读文档说明探索过程。实验据此阅读课程课件与官方 VecAdd 页面\cite{guide,ascend,bisheng}，在 \path{lab0/prep/03_mlir/ascend_vecadd.mlir} 保留注明来源的参考程序。

官方示例处理 16 个 \texttt{i16} 元素。参数位于 GM 地址空间，两个 \texttt{hivm.hir.load} 将输入搬入 UB，\texttt{hivm.hir.vadd} 在片上完成加法，再由 \texttt{hivm.hir.store} 写回 GM。\texttt{hacc.entry} 与 DEVICE 属性表达设备入口。该程序已经包含硬件相关操作，不能把它称为未经降低的通用高层算子。

\begin{table}[H]
\centering\small
\caption{课程材料中的 AscendNPU IR 分工}\label{tab:ascend}
\begin{tabularx}{\textwidth}{lXX}
\toprule
表示 & 主要表达内容 & VecAdd 探索中的关注点\\
\midrule
HFusion & 较高层算子及融合相关信息 & 在结构尚保留时进行变换与融合\\
HIVM & 设备操作、地址空间、搬运及同步约束 & GM/UB、load/vadd/store\\
HACC & 异构执行、设备入口与函数属性 & entry、DEVICE 与主机／设备协作\\
\bottomrule
\end{tabularx}
\end{table}
这些方言承载不同职责，可在表示中协作；HACC 不是简单排在 HIVM 之后、替换全部算子的单一流水级。逐层降低让高层优化保留所需语义，再逐步落到存储、指令和硬件资源约束。

\subsection{从官方示例推导存储与同步约束}
为使调研落到具体程序，沿官方例子的两次 load、一次 vadd、一次 store 追踪数据依赖，并与架构文档中的核内存储和流水映射职责对照\cite{ascendarch}。三个 UB 缓冲区的逻辑载荷各为 $16\times2=32$ 字节，合计 96 字节；GM 的逻辑读写量为 $2\times32+32=96$ 字节。这只是按元素数推导的有效载荷，实际分配可能包含对齐、临时空间或复用，不能冒充内存规划 pass 的实测结果。

\Needspace{9\baselineskip}
图\ref{fig:ascenddeps}给出需要保持的偏序：向量加必须等两个输入搬入 UB 后才可执行，写回必须等待加法结果。该图表达数据依赖，不指定真实事件号或同步指令；专用编译器还需将依赖映射到硬件流水。
\begin{figure}[H]
\centering\small
\begin{tikzpicture}[>=stealth,every node/.style={font=\small},
  box/.style={draw=black!45,rounded corners=2pt,align=center,minimum height=0.85cm}]
\node[box] (a) at (0,0.6) {GM 输入 A\\load 到 UB};
\node[box] (b) at (0,-0.6) {GM 输入 B\\load 到 UB};
\node[box] (add) at (3.4,0) {UB 内 vadd\\16 个 i16 相加};
\node[box] (out) at (6.7,0) {store 到 GM\\结果 C};
\draw[->] (a.east) -- (add.west);
\draw[->] (b.east) -- (add.west);
\draw[->] (add.east) -- (out.west);
\end{tikzpicture}
\caption{由官方 VecAdd 源码推导的数据依赖；非设备实测时间线}\label{fig:ascenddeps}
\end{figure}
这个小样例只有一个向量加，没有相邻算子可融合，也不含 Cube 运算，因而不能用它声称已验证算子融合收益或 Cube/Vector 混合核拆分。它适合观察地址空间、搬运、内存规划和同步约束何时进入 IR。课件中的地址偏移、同步伪码与 \texttt{llvm.npu.*} 名称用于解释原理；本报告没有将其作为特定版本的真实 pass 输出。

官方设备端编译方式为：
\par\noindent\begin{minipage}{\linewidth}
\begin{lstlisting}[numbers=none]
bishengir-compile add.mlir -enable-hivm-compile -o kernel.o
\end{lstlisting}
\end{minipage}\par
随后由 CANN host 程序加载或注册设备二进制、启动 kernel 并取回结果。本次环境的 PATH 中未找到 \texttt{bishengir-compile}、\texttt{bishengir-opt}、\texttt{hivmc} 或 \texttt{npu-smi}，默认 CANN 安装目录也不存在。实际用 \texttt{mlir-opt-15} 读取参考程序时，工具在 \texttt{hivm.address\_space} 处报告 \texttt{unregistered dialect} 并返回 1，诊断保存在 \path{lab0/report/evidence/ascend-probe.json}。允许未注册方言只能绕过相应解析限制，不能补上缺失的验证器和 lowering pass；因此未生成设备二进制，也未采集专用 pass 转储。

阅读官方 host 代码时还发现，其比较循环打印 \texttt{Expect} 和 \texttt{Result}，随后直接打印成功提示；循环本身未因数值不等而报错。因此，未来上板验证应检查全部 16 个元素相等并在失败时返回非零值，同时检查运行时调用结果，不能把示例中的成功字符串单独当作正确性证据。通用 MLIR 的 C 驱动已经采用逐元素检查并返回失败状态。

\subsection{比较与探索边界}
通用实验明确展示了“结构化循环 $\rightarrow$ 显式 CFG $\rightarrow$ LLVM 方言 $\rightarrow$ LLVM IR”的实测过程。Ascend 文档则补充了 CPU 样例中不显式出现的片上存储、数据搬运和设备入口。前者验证逐层转换后的数值结果，后者用于理解专用方言的设计。两部分共同说明多级表示的价值，但通用实验通过并不能证明 Ascend 参考程序在 NPU 上通过；设备专用调度、同步和性能仍需相应工具链与硬件验证。

\Needspace{15\baselineskip}
\section{结论}
本实验通过阶段产物建立了源程序与目标程序之间可检查的联系。预处理改变编译输入，前端组织并检查程序语义，中端将控制流及数据流显式化并优化，后端结合目标指令和调用约定生成汇编，汇编器与链接器再完成机器码、符号与地址的组织。

等价编程的 170 次比对和编译流程的 85 次检查均通过。短路副作用、二维数组寻址、递归调用、负数除余和浮点舍入，是低层实现中需要维护的语义；SSA 提升与 O2 对照则表明，正确性、指令结构、文件体积和性能是不同问题。MLIR 四级转储及 12 个输出元素验证说明逐层转换如何保留计算；Ascend 调研则进一步把向量加映射到地址空间、搬运和同步依赖。后续课程编译器可复用这些样例与运行环境，逐步实现语法分析、类型检查、SSA、优化及自有 RISC-V 后端。

\appendix
\ctexset{section={name={附录,},number=\Alph{section}}}
\renewcommand{\sectionmark}[1]{\markboth{附录\Alph{section}\quad #1}{}}
\section{复现与结果文件}
在 Ubuntu／WSL 安装表\ref{tab:tools}中的工具，并将课程运行库放入 \path{lab0/runtime/lib.tar.gz} 后，在仓库根目录执行：
\par\noindent\begin{minipage}{\linewidth}
\begin{lstlisting}[numbers=none]
make prep       # compilation stages, 170 equivalence checks, MLIR
make test       # current compiler lexer regression tests
make report     # export PDF from lab0/report/preparation.tex
make clean      # remove rebuildable intermediates
\end{lstlisting}
\end{minipage}\par
\path{lab0/report/evidence/results.json} 保存执行时间、工具版本、源码及运行库 SHA-256、命令、输入和实际输出；\path{lab0/report/evidence/observations.txt} 保存报告使用的 GIMPLE、IR、汇编、诊断及 MLIR 各层输出。\path{lab0/build/} 中的完整日志与二进制可重新生成，清理不会删除上述报告依据。词法分析器测试仅验证当前工程状态，不计入本次预备实验的 170 次等价比对。

\clearpage
\begin{thebibliography}{10}
\addcontentsline{toc}{section}{参考文献}
\bibitem{assignment} 课程资料. 预备工作：了解你的编译器，上机作业要求与作业完成要求.
\bibitem{sysy} 课程资料. SysY2022 语言定义 V1；SysY 文法补充说明.
\bibitem{runtime} 课程资料. SysY2022 运行时库 V1；lib.tar.gz.
\bibitem{guide} 课程资料. 编译器开发环境部署；了解编译器与 LLVM IR 编程；AscendNPU-IR 课件.
\bibitem{gcc} GNU Project. GCC 12.3 Manual: Options Controlling the Kind of Output. \url{https://gcc.gnu.org/onlinedocs/gcc-12.3.0/gcc/Overall-Options.html}.
\bibitem{llvm} LLVM Project. LLVM Language Reference Manual. \url{https://llvm.org/docs/LangRef.html}.
\bibitem{riscv} RISC-V International. RISC-V ABIs Specification. \url{https://riscv-non-isa.github.io/riscv-elf-psabi-doc/}.
\bibitem{mlir} LLVM Project. MLIR Documentation. \url{https://mlir.llvm.org/docs/}.
\bibitem{ascend} Huawei. AscendNPU IR：编译与执行示例. \url{https://ascendnpu-ir.gitcode.com/zh_cn/sources/introduction/quick_start/examples_zh.html}. 阅读日期：2026年10月10日.
\bibitem{ascendarch} Huawei. AscendNPU IR：架构设计. \url{https://ascendnpu-ir.gitcode.com/zh_cn/sources/introduction/architecture_zh.html}. 阅读日期：2026年10月10日.
\bibitem{bisheng} Huawei. 毕昇编译器与 CANN. \url{https://www.hiascend.com/cann/bisheng}.
\end{thebibliography}
\end{document}
